\documentclass[10pt]{article}

% ---------- Document Identity (update these only) ----------
\newcommand{\DocStage}{}
\newcommand{\DocTitle}{Your Road to Tier One \textbf{SOF}}
\newcommand{\ProgramName}{182-Week Civilian-to-Selection Readiness Pipeline}
\newcommand{\ProgramSubtitle}{}

% ---------- Page ----------
\usepackage[legalpaper,left=0.5in,right=0.5in,top=0.6in,bottom=0.45in]{geometry}

% ---------- Typography ----------
\usepackage[T1]{fontenc}
\usepackage[utf8]{inputenc}
\usepackage{libertinus}
\usepackage{microtype}
\usepackage{parskip}
\usepackage{ragged2e}
\usepackage{textcomp}
\AtBeginDocument{\color{black}}
\AtBeginDocument{\pagecolor{white}}
\AtBeginDocument{\justifying}

% ---------- Landscape pages ----------
\usepackage{pdflscape}

% ---------- Color ----------
\usepackage[table]{xcolor}
\definecolor{StageBlue}{HTML}{1F4E79}
\definecolor{StageBlueDark}{HTML}{163A5A}
\definecolor{LightGray}{HTML}{F2F4F7}
\definecolor{MidGray}{HTML}{D9DEE5}
\definecolor{RuleGray}{HTML}{C7CED8}
\definecolor{HeaderInk}{HTML}{000000}
\definecolor{Ink}{HTML}{000000}

% ---------- Utility ----------
\usepackage{enumitem}
\sloppy
\setlength{\emergencystretch}{2em}

% ---------- Boxes / UI ----------
\usepackage[most]{tcolorbox}

\tcbset{
  charterTitle/.style={
    enhanced,
    colback=StageBlue!4,
    colframe=StageBlue,
    boxrule=1.25pt,
    arc=3pt,
    left=16pt,right=16pt,top=14pt,bottom=14pt,
    borderline north={3.2pt}{0pt}{StageBlue},
    borderline west={4pt}{0pt}{StageBlue},
    borderline south={1.25pt}{0pt}{StageBlue},
    drop shadow={black!25!white},
  },
  charterRule/.style={
    enhanced,
    colback=LightGray,
    colframe=RuleGray,
    boxrule=0.85pt,
    arc=3pt,
    left=12pt,right=12pt,top=10pt,bottom=10pt,
    borderline west={3pt}{0pt}{StageBlue},
  },
  charterBadge/.style={
    enhanced,
    colback=StageBlue,
    coltext=white,
    colframe=StageBlueDark,
    boxrule=0.6pt,
    arc=2pt,
    left=6pt,right=6pt,top=3pt,bottom=3pt,
  }
}
\newtcbox{\Badge}{nobeforeafter,charterBadge}

% ---------- Lists ----------
\setlist[itemize]{leftmargin=1.5em, itemsep=0.25em, topsep=0.25em}
\setlist[enumerate]{leftmargin=1.8em, itemsep=0.25em, topsep=0.25em}

% ---------- TOC ----------
\usepackage{tocloft}
\setcounter{tocdepth}{3}
\setlength{\cftbeforesecskip}{3pt}
\setlength{\cftbeforesubsecskip}{2pt}
\setlength{\cftbeforesubsubsecskip}{0pt}
\renewcommand{\cftsecfont}{\bfseries}
\renewcommand{\cftsecpagefont}{\bfseries}
\renewcommand{\cftsubsecfont}{\normalfont}
\renewcommand{\cftsubsecpagefont}{\normalfont}
\renewcommand{\cftsubsubsecfont}{\normalfont}
\renewcommand{\cftsubsubsecpagefont}{\normalfont}
\setlength{\cftsecindent}{0pt}
\setlength{\cftsubsecindent}{1.25em}
\setlength{\cftsubsubsecindent}{2.8em}
\setlength{\cftsecnumwidth}{2.1em}
\setlength{\cftsubsecnumwidth}{2.8em}
\setlength{\cftsubsubsecnumwidth}{3.6em}

% Remove the built-in TOC heading so only the custom heading is shown
\makeatletter
\renewcommand{\tableofcontents}{\@starttoc{toc}}
\makeatother

% ---------- Header/Footer: page number only ----------
\usepackage{fancyhdr}
\pagestyle{fancy}
\fancyhf{}
\renewcommand{\headrulewidth}{0pt}
\renewcommand{\footrulewidth}{0pt}
\fancyfoot[C]{\thepage}

% ---------- Section formatting ----------
\usepackage{titlesec}

\titleformat{\section}
  {\Large\bfseries\color{Ink}}
  {\thesection}{0.6em}{}
  [\vspace{0.25em}\textcolor{StageBlue}{\rule{\linewidth}{0.8pt}}]

\titleformat{\subsection}
  {\large\bfseries\color{Ink}}
  {\thesubsection}{0.6em}{}

\titleformat{\subsubsection}
  {\normalsize\bfseries\color{Ink}}
  {\thesubsubsection}{0.6em}{}

\titlespacing*{\section}{0pt}{0.95em}{0.35em}
\titlespacing*{\subsection}{0pt}{0.65em}{0.25em}
\titlespacing*{\subsubsection}{0pt}{0.55em}{0.2em}

% ---------- Table engine ----------
\usepackage{tabularray}
\UseTblrLibrary{booktabs}

% ---------- tabularray: GLOBAL table treatment ----------
\SetTblrInner{
  cells = {fg=black, bg=white, valign=t},
}

% ---------- Header helpers ----------
\newcommand{\Hdr}[1]{\shortstack[c]{\strut #1\strut}}

% ---------- Lines ----------
\newcommand{\TitleLine}{\textcolor{StageBlue}{\rule{0.98\linewidth}{0.95pt}}}
\newcommand{\SubTitleLine}{\textcolor{RuleGray}{\rule{0.98\linewidth}{0.65pt}}}

% ---------- Continued on next page ----------
\newcommand{\ContinuedOnNextPageInline}{%
  \par
  \vfill
  \noindent\textcolor{RuleGray}{\rule{\linewidth}{1pt}}\vspace{-2pt}
  \noindent\hfill\textit{Continued on next page}\par\vspace{-10pt}
  \noindent\textcolor{RuleGray}{\rule{\linewidth}{1pt}}
  \clearpage
}

% ---------- PDF hyperlinks ----------
\usepackage[hidelinks]{hyperref}
\hypersetup{
  colorlinks=false,
  pdfborder={0 0 0},
  linktoc=all,
  pdftitle={\DocTitle},
  pdfauthor={\ProgramName}
}

% =========================================================
\begin{document}

\begin{tcolorbox}
\begin{center}
Stage 10 — Nutrition Operating System\\[0.35em]
\textbf{SOF} Physical Development Transformation Program\\[0.25em]
182-Week Civilian-to-Selection Readiness Pipeline\\[0.65em]
\TitleLine
\end{center}
\end{tcolorbox}

\vspace{0.35em}

\section*{Stage 10 — Nutrition Operating System}
\phantomsection
\addcontentsline{toc}{section}{Stage 10 — Nutrition Operating System}

\subsection*{Introduction}
\phantomsection
\addcontentsline{toc}{subsection}{Introduction}

\subsubsection*{1) Purpose}

Stage 10 defines the Nutrition Operating System for the 182-week pipeline. It provides a governed, operational nutrition framework that supports body composition change from the fixed start-state to the end-state envelope while protecting recovery, training performance, and the validity of gate evidence. Stage 10 is a control system and decision ladder; it does not override training governance, test protocols, phase purposes, or equipment timelines defined upstream.

Stage 10 exists to replace improvisation with auditable rules: stable inputs, controlled adjustments, conservative guardrails, and timing discipline aligned to protected test windows.

\subsubsection*{2) Scope}

\paragraph*{Inclusions}

\begin{itemize}
\item Stage 10 internal deliverable set (listed for completeness; not output here):
  \begin{itemize}
  \item 10.A Macro Ladder Phases 00–13
  \item 10.B Adjustment Rules
  \item 10.C Hydration and Electrolytes Operating System
  \item 10.D Troubleshooting Ladders
  \end{itemize}
\item Macro ladder by phase (10.A) defining a phase-appropriate operating posture that supports the training continuum without redefining it.
\item Adjustment rules (10.B) defining when and how changes are made (small-step, hold periods, trend-based) and how to avoid churn that contaminates validity.
\item Hydration/electrolytes operating system (10.C) defining numeric hydration/electrolyte targets and operational controls that preserve performance and safety without unsafe practices.
\item Troubleshooting ladders (10.D) defining common failure patterns (hunger spikes, adherence collapse, stalls, GI tolerance issues) and conservative corrective pathways.
\item Measurement interface anchored to Stage 2 and Stage 2.E:
  \begin{itemize}
  \item weekly bodyweight average,
  \item circumference suite,
  \item recovery markers,
  \item performance markers,
  \item validity tags and admissibility posture for decision use.
  \end{itemize}
\item Timing interface anchored to Stage 3:
  \begin{itemize}
  \item Deload/Test windows,
  \item reslot posture and no-stacking discipline,
  \item “no novelty near gates/test weeks” posture for diet variables.
  \end{itemize}
\end{itemize}

Required safety postures declared as binding for Stage 10:

\begin{itemize}
\item Weight-loss guardrail: average planned loss must not exceed 0.5\% bodyweight per week.
\item No planned refeeds or diet breaks.
\item Adjustments occur only via Stage 10.2 rules.
\item No novelty near gates/test weeks posture is enforced as a controlled-variable rule.
\item Red-flag symptom stop posture and escalation is preserved (non-medical referral posture).
\end{itemize}

Supplements posture in Stage 10:

\begin{itemize}
\item Supplements may be referenced only by category IDs 7.03–7.06 as optional supports.
\item 7.02 remains education-only and never required.
\item No supplement dosing guidance is provided in the Stage 10 Introductory Page.
\end{itemize}

\paragraph*{Exclusions}

\begin{itemize}
\item No training programming, workouts, or performance prescriptions.
\item No medical diagnosis or treatment guidance.
\item No prescription drug guidance.
\item No hormone dosing/sourcing/cycling/procurement guidance.
\item No meal plans, recipes, grocery lists, or product catalogs.
\item No supplement dosing; only category-level optional references are permitted.
\end{itemize}

\subsubsection*{3) How to Use}

\paragraph*{Coach Lead}

\begin{itemize}
\item Implement Stage 10 alongside training phases as a stability-first operating system:
  \begin{itemize}
  \item preserve controlled variables near protected test windows,
  \item avoid rapid diet changes that alter recovery and invalidate comparisons,
  \item enforce the 0.5\% bodyweight-per-week planned loss guardrail as a safety and performance control.
  \end{itemize}
\item Use Stage 2 metrics and Stage 2.E logging fields as the decision interface:
  \begin{itemize}
  \item decisions are trend-based, not single-day reactive.
  \end{itemize}
\item Apply Stage 10.B adjustment rules only at defined review points and with hold periods, especially around Deload/Test windows.
\item When conflicts with upstream constraints are detected (e.g., a nutrition change that would break “no novelty near gates”), flag Requires Stage 8 review rather than altering upstream doctrine.
\end{itemize}

\paragraph*{Trainee}

\begin{itemize}
\item Execute daily tracking and weekly review discipline using the prescribed measurement fields:
  \begin{itemize}
  \item record bodyweight per posture and compute weekly average,
  \item record circumference suite per cadence,
  \item record recovery markers (sleep, soreness, fatigue, pain flags) for readiness context,
  \item maintain honest logging; do not backfill by memory.
  \end{itemize}
\item Maintain stability:
  \begin{itemize}
  \item do not improvise changes week-to-week,
  \item do not introduce new diet variables during Deload/Test windows,
  \item follow the adjustment ladder rather than “making up” for bad days with extremes.
  \end{itemize}
\end{itemize}

\paragraph*{Test-week posture (validity protection)}

\begin{itemize}
\item Treat Deload+Test windows as protected from novelty:
  \begin{itemize}
  \item maintain stable nutrition inputs to avoid confounding performance evidence,
  \item if a deviation occurs, log it as a controlled-variable deviation and interpret test results conservatively.
  \end{itemize}
\item Do not use test weeks to introduce new restrictions, new hydration patterns, or novel supplement categories.
\end{itemize}

\subsubsection*{4) Inputs and Assumptions}

\paragraph*{Inputs}

\begin{itemize}
\item Stage 0 start-state and end-state envelope anchors:
  \begin{itemize}
  \item fixed start-state (age/height/weight/body fat, sedentary posture),
  \item end-state targets (bodyweight and body fat envelope, performance targets).
  \end{itemize}
\item Phase structure and durations from Stages 4–5 (Phase 00–13, fixed durations).
\item Measurement fields from Stage 2 and Stage 2.E:
  \begin{itemize}
  \item bodyweight weekly average,
  \item circumference suite,
  \item recovery markers,
  \item compliance markers,
  \item missing-data doctrine (record \textbf{MISSING}/\textbf{UNKNOWN}; no inference).
  \end{itemize}
\item Scheduling and reslot posture from Stage 3:
  \begin{itemize}
  \item Deload+Test windows and protected timing,
  \item adjustment rules for illness/regression/risk flags,
  \item no novelty near gates/test weeks.
  \end{itemize}
\item Tooling feasibility from Stage 7.13 minimal-path posture:
  \begin{itemize}
  \item basic container and hydration systems,
  \item cleaning and food safety posture,
  \item no reliance on specialty or expensive infrastructure for baseline execution.
  \end{itemize}
\end{itemize}

\paragraph*{Assumption Ledger only if needed}

\begin{itemize}
\item A tracking method is available (paper or digital), with offline fallback if needed.
\item Ability to weigh and measure weekly using stable tools and stable routines.
\end{itemize}

\subsubsection*{5) Interfaces}

\paragraph*{Upstream dependencies}

\begin{itemize}
\item Stage 0–5 constraints and phase sequencing govern training dose and gates; Stage 10 must support them without altering them.
\item Stage 2 measurement and Stage 2.E logging conventions provide the metrics and required fields used for nutrition decisions.
\item Stage 3 timing and reslot rules provide protected windows and no-novelty constraints that Stage 10 must respect.
\item Stage 7.13 provides the tooling layer for food prep and hydration system feasibility.
\end{itemize}

\paragraph*{Downstream consumers}

\begin{itemize}
\item Execution decisions and gate validity protection:
  \begin{itemize}
  \item Stage 10 ensures nutrition does not become a hidden confounder during gate evidence collection.
  \end{itemize}
\item Packaging and versioning updates:
  \begin{itemize}
  \item any changes to Stage 10 must be controlled and recorded with patch discipline and revalidation against Stage 3 protected windows and Stage 2 measurement posture.
  \end{itemize}
\end{itemize}

\paragraph*{Crosswalk hooks}

\begin{itemize}
\item Body composition measures:
  \begin{itemize}
  \item weekly bodyweight average,
  \item circumference suite,
  \item body fat estimate method recorded.
  \end{itemize}
\item Recovery markers:
  \begin{itemize}
  \item sleep hours/quality,
  \item soreness/fatigue,
  \item pain flags and foot flags.
  \end{itemize}
\item Deload/test week markers:
  \begin{itemize}
  \item protected windows and “no novelty” enforcement points.
  \end{itemize}
\item Change-log and patch discipline:
  \begin{itemize}
  \item patches recorded, impact level identified, and Stage 8 crosswalk impact review invoked when relevant.
  \end{itemize}
\end{itemize}

\subsubsection*{6) Gate / Acceptance Criteria for Stage 10}

Stage 10 is complete and deployable when:

\begin{itemize}
\item All four deliverables (10.1–10.4) are produced and internally consistent with Stages 0–9 constraints.
\item Numeric hydration/electrolyte targets exist in 10.3 (as required by Stage 10 scope) and are framed as operational targets without unsafe practices.
\item Adjustment rules (10.2) enforce:
  \begin{itemize}
  \item small-step changes,
  \item defined hold periods,
  \item trend-based decisions using Stage 2 metrics and Stage 2.E logging posture,
  \item the 0.5\% bodyweight-per-week planned loss guardrail.
  \end{itemize}
\item No planned refeeds or diet breaks posture is preserved.
\item No novelty near gates/test weeks posture is preserved and explicitly operationalized.
\item Feasibility under start-from-zero resources is preserved and does not require specialized equipment beyond Stage 7.13 minimal-path posture.
\item Any detected conflict with upstream doctrine is flagged as Requires Stage 8 review rather than silently altering upstream content.
\end{itemize}

\subsubsection*{7) Common Failure Modes + Controls}

\begin{itemize}
\item Over-aggressive restriction → enforce 0.5\% bodyweight-per-week planned loss guardrail; use maintenance reset posture via Stage 10.2 rather than escalating restriction.
\item Novelty near gate tests → enforce no novelty near gates/test weeks; stabilize inputs before protected windows; log deviations and interpret tests conservatively.
\item Treating supplements as required → prohibit; reference only optional category-level supports 7.03–7.06; 7.02 remains education-only and never required.
\item Poor data quality → enforce Stage 2.E field discipline; record \textbf{MISSING}/\textbf{UNKNOWN}; do not infer; do not make decisions from incomplete data.
\item Dehydration extremes → prohibit; Stage 10.3 must explicitly prevent water-cut practices and define safe operational rules; red-flag symptom stop posture applies.
\item Churn and weekly over-adjustment → enforce Stage 10.2 hold periods and review cadence; changes occur only by rule, not by mood.
\end{itemize}

\subsubsection*{8) Versioning Notes}

\begin{itemize}
\item Stage 10 is v1.0 aligned to program doctrine and the Stage 3 protected-window validity posture.
\item Changes require change control:
  \begin{itemize}
  \item patch records with impact level and revalidation scope,
  \item revalidation of interactions with test weeks and gate validity,
  \item conflicts flagged as Requires Stage 8 review rather than silent edits.
  \end{itemize}
\item Cosmetic corrections are permitted only when meaning is unchanged; substantive changes require logged patching and crosswalk impact assessment.
\end{itemize}

\end{document}